\section{Introduction}
Radio frequency (RF) fingerprinting attempts to identify a transmitter from physical characteristics of its received waveform. In the closed-set setting, each query is assigned to one of the transmitter identities represented in training. This task differs from detecting an unknown transmitter or deciding whether a packet is out of distribution. The observed waveform nevertheless combines transmitter imperfections with propagation and receiver effects. A successful classifier can therefore exploit patterns that fail when the receiver changes~\cite{wisig,shen}.

The operational distinction is important. A source-only model must identify a query using its trained parameters and the query waveform alone. A receiver-adaptation procedure is instead allowed to inspect a bounded set of unlabeled packets from the new receiver before classifying disjoint queries. A context-conditioning model can use those same packets without changing its learned weights. All three can be evaluated on an unseen-receiver protocol, but only the first is purely inductive domain generalization. Calling every method domain generalization obscures an information advantage that can be as important as architecture.

This study asks: among source-only generalization and unlabeled test-time receiver adaptation strategies, which methods improve transmitter recognition under unseen physical receiver shift when target information budgets are matched? A secondary question asks how that improvement changes with the size of the unlabeled support bank. Our unit of evaluation is the physical receiver rather than an independently resampled RF packet. Multiple seeds measure algorithmic variability within each receiver protocol.

The study emerged through a documented sequence of negative and attenuated results. An initial context study motivated a stricter disjoint-support evaluation. In that evaluation, attentive conditioning, denoted P2, did not provide a meaningful average improvement over independent empirical risk minimization (ERM), denoted P0. T3A~\cite{t3a}, an established classifier-adjustment procedure, performed better with the same unlabeled receiver information. We preserve those outcomes rather than redesign P2. The present manuscript consolidates frozen experiments and post-hoc audits; it reports no newly optimized model and no independent dataset replication.

Our contribution is a bounded empirical comparison and an auditable protocol, not a claim of algorithmic novelty for T3A or attention. The comparison separates source-only, unlabeled-support, and supervised-oracle information regimes; fixes support and query identities; retains every receiver and seed; and reports receiver-level effect heterogeneity, support dependence, probability calibration, and logged computational cost. It also exposes exclusions: BatchNorm-specific adaptation is incompatible with the frozen GroupNorm backbone, and a faithful Shen receiver-agnostic reproduction requires a different signal representation. These gaps limit any claim about the best adaptation method in general.

The principal finding is specific but useful: on this WiSig subset, a lightweight prototype adjustment improves unseen-receiver recognition more consistently than learned attentive context conditioning. The result justifies an independently frozen external or prospective replication. It does not establish real acquisition-episode realism, universal receiver robustness, or cross-dataset superiority.
